% File: unfaithful_cut_foundations.tex
%
% Main driver for the "Unfaithful Cut" foundations paper.
% Assumes section files live in ./sections and bibliography in ./bib.
%
% If you use \include instead of \input, do a clean build after renames
% (remove *.aux) to avoid stale include aux artifacts.

\documentclass[11pt]{article}

% ---------- Packages ----------
\usepackage{iftex}
\ifPDFTeX
  \usepackage[T1]{fontenc}
  \usepackage[utf8]{inputenc}
  \usepackage{lmodern}
\else
  \usepackage{fontspec}
  \defaultfontfeatures{Ligatures=TeX}
\fi
\usepackage{microtype}

\usepackage{amsmath,amssymb,amsthm,mathtools}
\usepackage{bm}

\usepackage{graphicx}
\usepackage{booktabs}
\usepackage{enumitem}

\usepackage{tikz}
\usetikzlibrary{arrows.meta,calc,positioning,decorations.pathmorphing}

\usepackage{hyperref}
\hypersetup{
  colorlinks=true,
  linkcolor=blue,
  citecolor=blue,
  urlcolor=blue,
  pdftitle={The Unfaithful Cut: A Structural Diagnostic for Algebraically Closed Reductions},
  pdfauthor={A. R. Wells},
}

\usepackage[margin=1in]{geometry}

% ---------- Bibliography ----------
% Uses biblatex (biber). If you prefer natbib/bibtex, say so and I’ll swap it.
\usepackage[
  backend=biber,
  style=numeric,
  sorting=nyt,
  maxbibnames=99
]{biblatex}
\addbibresource{bib/references.bib}

% ---------- Theorem Environments ----------
\theoremstyle{definition}
\newtheorem{definition}{Definition}[section]

\theoremstyle{plain}
\newtheorem{theorem}[definition]{Theorem}
\newtheorem{proposition}[definition]{Proposition}
\newtheorem{corollary}[definition]{Corollary}
\newtheorem{lemma}[definition]{Lemma}

\theoremstyle{remark}
\newtheorem{remark}[definition]{Remark}

% ---------- Convenience Macros ----------
\newcommand{\Hspace}{\mathcal{H}}
\newcommand{\Rspace}{\mathcal{R}}
\newcommand{\Iclass}{\mathcal{I}}
\newcommand{\Cont}{\mathrm{Cont}}

\newcommand{\SigOne}{\Sigma_{1}}
\newcommand{\SigTwo}{\Sigma_{2}}

% ---------- Title Block ----------
\title{\textbf{The Unfaithful Cut: A Structural Diagnostic for Algebraically Closed Reductions}}
\author{
  A.\ R.\ Wells\\
  \textit{Dual-Frame Research Group}\\
}
\date{January 6, 2026}

\begin{document}
\maketitle

% ---------- Abstract ----------
\input{sections/abstract.text}

% Optional: uncomment if you want a TOC
% \tableofcontents
% \bigskip

% ---------- Main Sections ----------
% These reflect your current directory listing.
% If you ran the rename script I provided, update the three inputs at the end accordingly.
\section{Introduction}

\subsection{Reduced descriptions and their limits}

Quantum theory is extraordinarily successful as a predictive framework, yet persistent foundational and operational difficulties arise whenever reduced descriptions are treated as sufficient representations of physical systems. Contextuality theorems prohibit noncontextual value assignments to quantum states; preparation contextuality shows that operationally equivalent states can encode inequivalent physical situations; non-Markovian dynamics demonstrate that instantaneous reduced states often fail to predict future behavior; and the measurement problem highlights the inability of unitary state evolution to account for definite events without additional structure.

These issues are typically addressed by enlarging the descriptive object. When states fail, one introduces environments, memory kernels, hidden variables, or multi-time process tensors. While often effective in practice, these strategies obscure a more basic and prior question: \emph{under what structural conditions can a reduced description support counterfactual prediction under interventions at all?}

\subsection{Operational reduction and the unfaithful cut}

Physical theories operate by reducing an underlying space of admissible physical histories to a smaller set of operationally accessible descriptions. Formally, this reduction is a many-to-one mapping from admissible histories to reduced objects. Such reductions are unavoidable: they define the interface between theory and experiment.

We argue that discarding information is not itself the problem. The problem arises when the reduction creates a \emph{degenerate representation}: dynamically distinct histories are identified at the reduced level, thereby discarding distinctions that are causally relevant for future interventions. In this case, the reduced description is not merely incomplete in an ontological sense; it is operationally insufficient for counterfactual reasoning.

We refer to this failure mode as \emph{the unfaithful cut}. Crucially, the failure is not that the reduced object evolves incorrectly, but that no evolution rule defined solely on that object can, in general, be sufficient to support counterfactual analysis—i.e., predictions of what would occur under alternative future interventions within a specified intervention class. In such cases, searching for improved master equations or more accurate dynamical maps constitutes a category error: the representational cut itself is inadequate for the task.

\subsection{From state insufficiency to structural diagnosis}

The unfaithful cut reframes several familiar problems. In open quantum systems, it clarifies why reduced density matrices fail when initial system--environment correlations are present. In non-Markovian dynamics, it explains why memory effects cannot always be captured by time-local or even multi-time algebraic descriptions. In quantum control, it accounts for the empirical observation that identical reduced states can respond differently to identical control pulses.

Our central claim is conditional and diagnostic: under physically standard assumptions about interventions, any reduced description that treats an algebraically closed object as sufficient—whether instantaneous or multi-time (Σ$_1$ descriptions)—fails to preserve admissible continuation structure whenever the reduction discards intervention-relevant distinctions. This failure is structural rather than dynamical and does not imply that Σ$_1$ descriptions are universally invalid.

We formalize the minimal representational structure required to restore counterfactual semantics as \emph{second-order constraint geometry} (Σ$_2$): a geometry over admissible reduced trajectories rather than over instantaneous reduced states. Memory effects, hysteresis, and arrow-like directedness then arise as consequences of projection and degeneracy, rather than as fundamental dynamical asymmetries.

\subsection{Scope and contributions}

This work makes three contributions. First, we define faithfulness with respect to interventions as a criterion for reduced descriptions and show how its failure diagnoses a broad class of foundational and operational limitations. Second, we establish a structural impossibility result identifying sufficient conditions under which algebraically closed (Σ$_1$) descriptions cannot support counterfactual prediction under a given intervention class. Third, we characterize the minimal representational information that any faithful description must encode, formalized here as Σ$_2$ geometry, thereby clarifying the structural origin of memory and directedness effects.

This work does not propose new dynamics, new ontologies, or a replacement formalism for quantum theory. It provides a structural diagnostic identifying when reduced descriptions are incapable—by construction—of supporting counterfactual reasoning under interventions, and it characterizes the minimal representational structure required to restore such reasoning.
\section{Formal Framework}
\label{sec:definitions}

\subsection{Admissible histories}

Let $\mathcal{H}$ denote the set of admissible histories of a physical system. Elements
$h \in \mathcal{H}$ represent complete realizations of the system subject to underlying
physical constraints. No preferred global temporal parameterization is assumed on
$\mathcal{H}$; however, histories may support partial orderings induced by experimental
interventions or boundary conditions.

Throughout, $\Sigma_1$ (``Sigma-one'') denotes algebraically closed reduced descriptions, while
$\Sigma_2$ (``Sigma-two'') denotes trajectory-level admissibility structures.

\subsection{Operational reduction}

An operational description is obtained via a reduction map
\begin{equation}
\Pi : \mathcal{H} \rightarrow \mathcal{R},
\end{equation}
where $\mathcal{R}$ is a space of reduced descriptions, such as density matrices, quantum
channels, or static multi-time process tensors. The reduction $\Pi$ is generally many-to-one
and represents the information retained by a given experimental or theoretical interface.

A reduction is said to be \emph{degenerate} if there exist distinct histories
$h_1 \neq h_2 \in \mathcal{H}$ such that
\begin{equation}
\Pi(h_1) = \Pi(h_2).
\end{equation}

\subsection{Algebraically closed reduced descriptions (\texorpdfstring{$\Sigma_1$}{Sigma-1})}

A reduced description $r \in \mathcal{R}$ is said to be \emph{algebraically closed} if all
admissible future predictions are determined by a fixed algebraic object together with
composition rules internal to a prescribed algebraic structure. Examples include
instantaneous quantum states with completely positive maps, quantum channels, and static
multi-time process tensors equipped with contraction rules.

We refer to any framework relying exclusively on algebraically closed reduced descriptions
as a $\Sigma_1$ framework.

\subsection{Interventions and history composition}

Let $\mathcal{I}$ denote a specified class of physical interventions. Interventions are
externally applied operations performed after the reduction, acting only on degrees of
freedom represented in $\mathcal{R}$. Examples include bounded Hamiltonian control pulses,
completely positive maps, and standard measurement-and-feedback operations. Interventions
that directly access discarded degrees of freedom are excluded.

We assume that admissible histories support a composition relation with interventions.
For $h,h' \in \mathcal{H}$ and $I \in \mathcal{I}$, we write
\begin{equation}
h \sim_I h'
\end{equation}
to mean that $h'$ is a possible history resulting from applying intervention $I$ to history
$h$. No determinism or uniqueness is assumed; $\sim_I$ is a relation, not a function. 

Measurements performed on the retained system $S$ that do not selectively post-condition on
environment outcomes are included as admissible interventions; operations requiring access to,
or conditioning on, the state of discarded degrees of freedom are excluded.

We do not require the intervention class $\mathcal{I}$ to be closed under composition or to form
a group. However, for multi-step counterfactual reasoning, one would typically assume that if
$I_1, I_2 \in \mathcal{I}$, then their sequential composition $I_2 \circ I_1$ is also admissible.
All results of this work apply whether or not such closure holds.

\subsection{Continuation sets}

For a reduced description $r \in \mathcal{R}$, define the \emph{continuation set} under
$\mathcal{I}$ as
\begin{equation}
\mathrm{Cont}(r;\mathcal{I})
=
\left\{
(I, r') \;\middle|\;
\exists\, h,h' \in \mathcal{H} \text{ such that }
\Pi(h)=r,\; h \sim_I h',\; \Pi(h')=r'
\right\}.
\end{equation}

The continuation set therefore records, for each admissible intervention, the set of
possible reduced outcomes compatible with at least one underlying history consistent
with $r$. This definition makes explicit that continuation structure is inherited from
$\mathcal{H}$ via the reduction $\Pi$, without assuming any reduced-level dynamics.

\subsection{Faithfulness with respect to interventions}

\begin{definition}[Faithfulness]
A reduction $\Pi$ is \emph{faithful with respect to} $\mathcal{I}$ if, for all histories
$h_1,h_2 \in \mathcal{H}$,
\begin{equation}
\Pi(h_1) = \Pi(h_2)
\;\Rightarrow\;
\mathrm{Cont}(\Pi(h_1);\mathcal{I})
=
\mathrm{Cont}(\Pi(h_2);\mathcal{I}).
\end{equation}
\end{definition}

Faithfulness requires that identical reduced descriptions support identical sets of
counterfactual continuations under the specified intervention class.

\subsection{Counterfactual prediction (operational sense)}

We say that a reduced description $r \in \mathcal{R}$ supports \emph{counterfactual prediction}
for an intervention class $\mathcal{I}$ if, for each $I \in \mathcal{I}$, the reduced theory
assigns a unique (or uniquely specified) set of admissible reduced outcomes
\begin{equation}
\widehat{\mathrm{Out}}(r;I) \subseteq \mathcal{R}
\end{equation}
that is invariant over the preimage $\Pi^{-1}(r)$; equivalently, if $\mathrm{Cont}(r;\mathcal{I})$
is well-defined independent of which history in $\Pi^{-1}(r)$ is realized. In an unfaithful cut,
identical reduced descriptions correspond to different counterfactual responses to the same $I$,
so no Σ$_1$ dynamics defined solely on $\mathcal{R}$ can determine the correct answer.

\subsection{The unfaithful cut}

A reduction $\Pi$ constitutes an \emph{unfaithful cut} if it is degenerate and fails to be
faithful with respect to $\mathcal{I}$. In this case, the reduced description is structurally
incapable of supporting counterfactual prediction under $\mathcal{I}$, regardless of the
dynamical laws imposed at the reduced level.
\section{Generic Unfaithfulness and a Minimal Example}

\subsection{A conditional structural-unfaithfulness result}

The faithfulness condition places a strong constraint on the reduction $\Pi$: whenever
distinct histories are identified at the reduced level, their future responses to all
admissible interventions must remain indistinguishable. This can hold only if the
discarded distinctions are dynamically irrelevant for the specified intervention class
$\mathcal{I}$.

Rather than appealing to a measure-theoretic notion of typicality, we state a conditional
\emph{structural diagnostic}: if a reduction discards degrees of freedom that can become
dynamically relevant under interventions that are standard in open-system practice
(including free evolution), then the reduction fails to be faithful. In this sense,
faithfulness requires special structural coincidences (e.g., absence of correlations)
that are unstable under perturbations of histories, interactions, or intervention
richness.

\begin{proposition}[Unfaithfulness from discarded correlations under free evolution]
\label{prop:unfaithful-from-correlations}
Assume a decomposition of the underlying system into retained and discarded degrees of
freedom (e.g., $S$ and $E$), and suppose:
\begin{enumerate}
  \item (\emph{Degenerate reduction}) There exist histories $h_1 \neq h_2$ such that
        $\Pi(h_1)=\Pi(h_2)=r$, but $h_1$ and $h_2$ differ in correlations between retained
        and discarded degrees of freedom.
  \item (\emph{Intervention richness}) The intervention class $\mathcal{I}$ contains at least
        one ``passive'' intervention $I_{\Delta t}$ corresponding to free evolution of the
        full retained+discarded system over a duration $\Delta t$ (no access to discarded
        degrees of freedom by the experimenter).
  \item (\emph{Correlation sensitivity}) For the above $h_1,h_2$, free evolution yields
        successor histories $h_1' , h_2'$ with distinct reduced outcomes:
        $\Pi(h_1') \neq \Pi(h_2')$.
\end{enumerate}
Then $\Pi$ is unfaithful with respect to $\mathcal{I}$.
\end{proposition}

\begin{remark}[Interpretive status]
This result is diagnostic rather than prohibitive: it identifies sufficient structural
conditions under which algebraically closed reduced descriptions cannot support
counterfactual prediction under the intervention class $\mathcal{I}$. It does not claim
that such descriptions must always be abandoned, nor that they fail outside these
conditions.
\end{remark}

\begin{remark}[Meaning of structural genericity]
Genericity here is understood \emph{structurally} rather than statistically:
it refers to the instability of faithfulness under perturbations to correlation
structure or intervention classes, not to the frequency with which approximate
$\Sigma_1$ descriptions succeed in controlled experimental regimes.
\end{remark}

\begin{remark}[Regimes of $\Sigma_1$ adequacy]
Faithful $\Sigma_1$ behavior persists in important and widely used regimes, including
factorized initial conditions, weak coupling or short-time limits, and restricted
intervention classes for which discarded degrees of freedom never become causally
relevant. The present result clarifies the boundaries of these regimes rather than
challenging their empirical success.
\end{remark}

\subsection{Minimal unfaithful example with explicit unitary dynamics}

Consider a system $S$ coupled to an environment $E$, with a joint unitary evolution
$U_{SE}(\Delta t)$ representing free evolution over a duration $\Delta t$. Let the reduction
$\Pi$ retain only the system state, so $\Pi(h)$ corresponds to $\rho_S(t_0)$ at an operational
cut time $t_0$.

Define two admissible histories $h_1,h_2 \in \mathcal{H}$ whose joint states at $t_0$ are
pure but correlated in distinct ways:
\begin{align}
|\psi_1\rangle_{SE}(t_0) &= \sum_i \alpha_i \, |s_i\rangle_S \, |e_i^{(1)}\rangle_E, \\
|\psi_2\rangle_{SE}(t_0) &= \sum_i \alpha_i \, |s_i\rangle_S \, |e_i^{(2)}\rangle_E,
\end{align}
where $\{|s_i\rangle\}$ is an orthonormal basis on $S$ and the environment families
$\{|e_i^{(1)}\rangle\}$ and $\{|e_i^{(2)}\rangle\}$ differ (capturing distinct correlations).

Assume both histories reduce to the same system density operator:
\begin{equation}
\rho_S(t_0)
=
\mathrm{Tr}_E\!\left[ |\psi_1\rangle\langle\psi_1| \right]
=
\mathrm{Tr}_E\!\left[ |\psi_2\rangle\langle\psi_2| \right].
\end{equation}
In particular, in the basis $\{|s_i\rangle\}$ this reduced state has the operator form
\begin{equation}
\rho_S(t_0) = \sum_i |\alpha_i|^2 \, |s_i\rangle\langle s_i|,
\end{equation}
which makes explicit that $\rho_S(t_0)$ depends on the $\alpha_i$ but not on the environment
vectors $|e_i^{(k)}\rangle$. Thus $\Pi(h_1)=\Pi(h_2)=\rho_S(t_0)$, so the reduction is degenerate.

Let the intervention class $\mathcal{I}$ include the passive intervention $I_{\Delta t}$
corresponding to free evolution by $U_{SE}(\Delta t)$. The evolved joint states are
\begin{equation}
|\psi_k\rangle_{SE}(t_0+\Delta t)
=
U_{SE}(\Delta t)\,|\psi_k\rangle_{SE}(t_0),
\quad k \in \{1,2\}.
\end{equation}
The corresponding reduced states are
\begin{equation}
\rho_S^{(k)}(t_0+\Delta t)
=
\mathrm{Tr}_E\!\left[ |\psi_k\rangle\langle\psi_k| (t_0+\Delta t) \right].
\end{equation}
For generic interaction unitaries $U_{SE}(\Delta t)$ and distinct correlation structure
in the environment families, one obtains
\begin{equation}
\rho_S^{(1)}(t_0+\Delta t) \neq \rho_S^{(2)}(t_0+\Delta t),
\end{equation}
because $\rho_S(t_0)$ does not encode which environment correlation pattern is present.

This is a \textbf{predictive} failure: an experimenter who prepares $\rho_S(t_0)$ and knows only
the $\Sigma_1$ description cannot predict the outcome of applying $I_{\Delta t}$. A $\Sigma_1$ framework
would prescribe a unique reduced dynamical map $\mathcal{E}_{\Delta t}$ such that
\begin{equation}
\rho_S(t_0+\Delta t) = \mathcal{E}_{\Delta t}\!\left[\rho_S(t_0)\right],
\end{equation}
but in the presence of the unfaithful cut the same input $\rho_S(t_0)$ yields multiple distinct
outputs depending on hidden correlation structure. No choice of $\mathcal{E}_{\Delta t}$ can be
correct for both histories simultaneously without additional structure beyond the $\Sigma_1$ object.

\subsection{Classical Analog in Stochastic Control and Causal Inference}
\label{sec:classical-analog}

This subsection provides a non-quantum instantiation of the unfaithful cut.
It introduces no new claims for classical control or causal inference; it
serves only to demonstrate that unfaithful cuts arise whenever intervention
semantics are imposed on degenerate reductions, independent of quantum
structure.

\paragraph{Hidden-state control model.}
Let $X_t \in \mathcal{X}$ denote an unobserved state, $Y_t \in \mathcal{Y}$
an observed quantity, and $U_t \in \mathcal{U}$ a controllable intervention.
Assume a controlled stochastic evolution
\begin{equation}
  X_{t+1} = f(X_t, U_t, \xi_t),
  \qquad
  Y_t = g(X_t),
  \label{eq:classical-hidden-state}
\end{equation}
where $\{\xi_t\}$ is exogenous noise. A \emph{history} $h \in \mathcal{H}$
may be taken as a realization of $(X_0, Y_0, U_0, X_1, Y_1, U_1, \dots)$,
or more abstractly as an element of whatever admissible trajectory space
the model induces.

Define a reduction $\Pi:\mathcal{H}\to\mathcal{R}$ by retaining only the
current observation at a designated reduction event:
\begin{equation}
  \Pi(h) \coloneqq Y_t.
  \label{eq:classical-reduction}
\end{equation}
This is the classical analogue of working with an instantaneous reduced
descriptor and treating it as sufficient.

\paragraph{Interventions and continuation sets.}
Let the intervention class $\mathcal{I}$ include one-step control actions
$I_u$ that set $U_t=u$ at the next step and otherwise act only through the
retained degrees of freedom. Using the intervention-induced relation
$h \sim_{I} h'$ introduced in Section~\ref{sec:definitions}, define
continuation outcomes in the reduced space by
\begin{equation}
  \mathrm{Cont}(y;\mathcal{I})
  \;\coloneqq\;
  \Bigl\{ (I, y') \,:\,
  \exists\, h,h' \in \Pi^{-1}(y)\ \text{with}\ h \sim_{I} h'
  \ \wedge\ y'=\Pi(h') \Bigr\}.
  \label{eq:classical-cont}
\end{equation}
Faithfulness requires that if $\Pi(h_1)=\Pi(h_2)=y$, then
$\mathrm{Cont}(\Pi(h_1);\mathcal{I})=\mathrm{Cont}(\Pi(h_2);\mathcal{I})$.

\paragraph{An unfaithful cut from hidden confounding.}
Suppose there exist two admissible histories $h_1,h_2\in\mathcal{H}$ such that
\begin{equation}
  \Pi(h_1)=\Pi(h_2)=y,
  \qquad
  X_t^{(1)} \neq X_t^{(2)}.
  \label{eq:classical-same-observation}
\end{equation}
If the dynamics are correlation-sensitive in the sense that there exists a
control action $u$ for which
\begin{equation}
  \Pi(h_1') \neq \Pi(h_2')
  \quad \text{for some}\quad
  h_1 \sim_{I_u} h_1',\;\; h_2 \sim_{I_u} h_2',
  \label{eq:classical-diverge}
\end{equation}
then the reduction \eqref{eq:classical-reduction} is unfaithful: the same
reduced object $y$ supports distinct continuation structure under the same
intervention class $\mathcal{I}$.

This failure is the control-theoretic analogue of non-identifiability in causal
models: distinct latent states compatible with the same observation yield
different outcomes under identical do-interventions, as in Simpson's paradox driven by hidden confounding.

This is a predictive failure. Conditioning only on $Y_t=y$ does not suffice
to determine the outcome of applying $U_t=u$.

\paragraph{Concrete two-latent example.}
For intuition, take $\mathcal{X}=\{A,B\}$ and $\mathcal{Y}=\{0,1\}$ with a
degenerate observation map $g(A)=g(B)=0$. Let $\mathcal{U}=\{0,1\}$ and define
\begin{equation}
  \mathbb{P}(Y_{t+1}=1 \mid X_t=A, U_t=1)=1,
  \qquad
  \mathbb{P}(Y_{t+1}=1 \mid X_t=B, U_t=1)=0,
  \label{eq:classical-two-latent}
\end{equation}
(with any convenient specification for $U_t=0$). Then $Y_t=0$ is compatible
with both $X_t=A$ and $X_t=B$, but applying the same intervention $U_t=1$
yields deterministically distinct outcomes for $Y_{t+1}$.

\paragraph{$\Sigma_2$ resolution (classical form).}
As in the quantum case, the $\Sigma_2$ description replaces a point-valued
reduction by a fiber-aware object. The fiber at $y$ is
\begin{equation}
  \mathcal{F}(y) \coloneqq \Pi^{-1}(y),
  \label{eq:classical-fiber}
\end{equation}
equipped with the intervention-indexed admissibility relations inherited
from $\sim_I$. Interventions generically induce set-valued updates, reflecting
the fact that different admissible completions of the same reduced descriptor
support different future continuations.

\subsection{Addressing ``just add the environment''}

A natural response to the preceding example is to enlarge the reduced description to include
the environment, replacing $\rho_S$ by $\rho_{SE}$. This move can restore faithfulness for
that particular choice of ``system,'' but it does not eliminate the phenomenon: it merely
pushes the operational cut outward. Any finite retained subsystem can be correlated with
additional external degrees of freedom, and if the intervention class includes free evolution
of the larger closed dynamics, the same argument reappears.

\begin{corollary}[Pushing the cut outward does not remove unfaithfulness]
\label{cor:push-cut}
Let $\Pi$ describe any finite retained subsystem. If histories in $\mathcal{H}$ admit
correlations between the retained subsystem and discarded degrees of freedom, and if
$\mathcal{I}$ includes passive free evolution of the total system, then $\Pi$ is unfaithful
unless those correlations are excluded (or rendered dynamically irrelevant) by assumption.
\end{corollary}

Corollary~\ref{cor:push-cut} clarifies the structural content of the unfaithful cut: the issue
is not that one picked the ``wrong'' subsystem, but that any many-to-one operational reduction
that discards potentially relevant correlations will fail under sufficiently rich intervention
classes.
\section{\texorpdfstring{$\Sigma_2$}{Sigma-2} Geometry as the Minimal Faithfulness Structure}
\label{sec:sigma2}

Sections 2--3 formalize the unfaithful cut and establish that $\Sigma_1$ descriptions
can fail under physically standard conditions when a reduction discards intervention-relevant
distinctions. We now define a $\Sigma_2$ description and show, in a precise sense, that it
\emph{explicitly represents the minimal information any faithful reduced description must encode}
to support counterfactual prediction under a specified intervention class $\mathcal{I}$,
without collapsing back to the trivial solution of retaining all of $\mathcal{H}$.

\subsection{The \texorpdfstring{$\Sigma_2$}{Sigma-2} object: fibers with admissibility structure}

Let $\Pi:\mathcal{H}\to\mathcal{R}$ be a fixed reduction and $\mathcal{I}$ an intervention class.
For each $r\in\mathcal{R}$, define the \emph{fiber} (preimage)
\begin{equation}
\mathcal{F}(r) := \Pi^{-1}(r) \subseteq \mathcal{H}.
\end{equation}

Define the \emph{intervention successor relation on histories} as $\sim_I \subseteq \mathcal{H}\times \mathcal{H}$
(Section 2). This induces, for each $r$ and each $I\in\mathcal{I}$, a \emph{fiber-to-outcome relation}:
\begin{equation}
\mathcal{A}_I(r) \subseteq \mathcal{F}(r)\times\mathcal{R},
\qquad
(h, r')\in\mathcal{A}_I(r)
\iff
\exists h' \in \mathcal{H}: h \sim_I h' \;\wedge\; \Pi(h')=r'.
\end{equation}

Intuitively, $\mathcal{A}_I(r)$ records which reduced outcomes $r'$ are reachable under $I$,
\emph{indexed by which fiber element $h\in\Pi^{-1}(r)$ is realized}.

\begin{definition}[Σ$_2$ description]
A $\Sigma_2$ description of a reduced state $r\in\mathcal{R}$ (relative to $\Pi$ and $\mathcal{I}$)
is the structured object
\begin{equation}
\Sigma_2(r) := \bigl(r, \mathcal{F}(r), \{\mathcal{A}_I(r)\}_{I\in\mathcal{I}}\bigr),
\end{equation}
consisting of the base point $r$, its fiber $\mathcal{F}(r)$, and the family of admissibility
relations $\{\mathcal{A}_I(r)\}$ describing counterfactual continuation structure under $\mathcal{I}$.
\end{definition}

This definition is deliberately minimal in \emph{representational} scope: it adds no dynamics on
$\mathcal{R}$ and no privileged time parameter. It supplements $r$ by the least additional
structure required to make intervention-indexed continuation structure explicit.

\subsection{Composition under interventions}

A $\Sigma_2$ description must be operationally usable: given $\Sigma_2(r)$ and an intervention $I$,
we must be able to compute the resulting $\Sigma_2$ object(s). Define the \emph{post-intervention fiber}
at an outcome $r'\in\mathcal{R}$ by
\begin{equation}
\mathcal{F}_I(r \to r') := \left\{ h' \in \mathcal{H} \;\middle|\;
\exists h \in \mathcal{F}(r)\text{ with } h \sim_I h' \text{ and } \Pi(h')=r'
\right\}.
\end{equation}

Note that $\mathcal{F}_I(r \to r')$ is a subset of the target fiber:
\begin{equation}
\mathcal{F}_I(r \to r') \subseteq \mathcal{F}(r') = \Pi^{-1}(r'),
\end{equation}
consisting only of those fiber elements at $r'$ reachable via the intervention path from some
history in $\mathcal{F}(r)$. In this sense, interventions refine fibers rather than selecting
single histories.

The $\Sigma_2$ update induced by $I$ is set-valued:
\begin{equation}
\Sigma_2(r) \xRightarrow[]{I}
\left\{
\Sigma_2(r') \;:\; \exists\, h\in\mathcal{F}(r)\ \exists\,h'\in\mathcal{H}\ (h\sim_I h' \wedge \Pi(h')=r')
\right\}.
\end{equation}

For instance, if $\mathcal{F}(r)$ contains histories $\{h_1,h_2\}$ and an intervention $I$ yields
$h_1\sim_I h_1'$ with $\Pi(h_1')=r_1'$ while $h_2\sim_I h_2'$ with $\Pi(h_2')=r_2'$ and
$r_1'\neq r_2'$, then the update produces two distinct $\Sigma_2$ outcomes
$\{\Sigma_2(r_1'),\Sigma_2(r_2')\}$.

This rule is precisely what $\Sigma_1$ frameworks cannot supply in the unfaithful case: Σ$_1$ attempts
to assign a single reduced evolution rule $r \mapsto r'$ independent of which $h\in\Pi^{-1}(r)$
is realized, while $\Sigma_2$ makes this dependence explicit.

\subsection{Faithfulness restored}

We now show that $\Sigma_2$ restores counterfactual prediction relative to $\mathcal{I}$ by construction.

\begin{proposition}[$\Sigma_2$ determines continuation structure]
\label{prop:sigma2-determines-cont}
For each $r\in\mathcal{R}$, the continuation set $\mathrm{Cont}(r;\mathcal{I})$ is determined by
$\Sigma_2(r)$ via
\begin{equation}
\mathrm{Cont}(r;\mathcal{I})
=
\left\{
(I,r') \;\middle|\; \exists h\in\mathcal{F}(r)\ \text{such that}\ (h,r')\in\mathcal{A}_I(r)
\right\}.
\end{equation}
\end{proposition}

\begin{proof}
Immediate from the definitions of $\mathcal{F}(r)$ and $\mathcal{A}_I(r)$.
\end{proof}

\begin{corollary}[Counterfactual prediction becomes well-defined]
\label{cor:sigma2-counterfactual}
If a reduced theory represents preparation outcomes by $\Sigma_2$ objects, then counterfactual
prediction under $\mathcal{I}$ is well-defined in the sense that the theory can represent,
for each $I\in\mathcal{I}$, the full set of admissible reduced outcomes without ambiguity.
\end{corollary}

The point is not that $\Sigma_2$ collapses multiple outcomes to one; rather, it represents the
\emph{set-valued} nature of counterfactual outcomes induced by degeneracy in $\Pi$.

\subsection{Representational minimality of \texorpdfstring{$\Sigma_2$}{Sigma-2} (continuation-sufficient structure)}

We formalize ``minimal'' as \emph{representational}: $\Sigma_2$ makes explicit the coarsest
augmentation of $r$ that suffices to recover the continuation set for all interventions in
$\mathcal{I}$. Any faithful completion must encode at least the continuation structure, hence
at least the information content of $\Sigma_2$. This is a representational minimality claim,
not a computational one.

Let $\mathfrak{D}$ be any candidate augmented description space and $K:\mathcal{H}\to\mathfrak{D}$
a refined reduction (a ``completion'') such that $K$ is faithful with respect to $\mathcal{I}$
in the same sense as Definition 2.1 (with $\Pi$ replaced by $K$). Suppose further that
$\Pi$ factors through $K$, i.e.,
\begin{equation}
\Pi = q \circ K
\end{equation}
for some surjection $q:\mathfrak{D}\to\mathcal{R}$, meaning $K$ refines $\Pi$.

\begin{theorem}[Representational minimality of $\Sigma_2$ as continuation-sufficient structure]
\label{thm:minimal-sigma2}
Let $K:\mathcal{H}\to\mathfrak{D}$ be any completion refining $\Pi$ and faithful with respect
to $\mathcal{I}$. Then there exists a map
\begin{equation}
\phi:\mathfrak{D}\to \Sigma_2(\mathcal{R})
\end{equation}
such that for all $h\in\mathcal{H}$,
\begin{equation}
\phi(K(h)) = \Sigma_2(\Pi(h)),
\end{equation}
and such that $\phi(d)$ determines $\mathrm{Cont}(q(d);\mathcal{I})$.
In particular, any faithful completion must carry at least enough information to reconstruct
the $\Sigma_2$ continuation structure.
\end{theorem}

\begin{proof}[Proof sketch]
Because $K$ is faithful, for any two histories $h_1,h_2\in K^{-1}(d)$ we have identical continuation
structure under $\mathcal{I}$ at the $K$-level. Since $\Pi=q\circ K$, both histories map to the same
basepoint $r=q(d)\in\mathcal{R}$, and their induced continuation structure at the $\Pi$-level must
therefore coincide. Hence the fiber $K^{-1}(d)$ can be embedded into a well-defined subset of
$\Pi^{-1}(r)$ on which the continuation relations induced by $\sim_I$ are consistent. Define $\phi(d)$
to be the $\Sigma_2$ object at $r$ whose fiber is this embedded subset and whose admissibility relations
are induced by $\sim_I$ restricted to it. This construction is well-defined precisely because
faithfulness of $K$ ensures invariance of continuation structure over $K^{-1}(d)$.
\end{proof}

Theorem~\ref{thm:minimal-sigma2} establishes that $\Sigma_2$ encodes the \emph{representationally minimal} information
required for faithfulness. This does not imply computational minimality: in worst cases,
$\Sigma_2$ may be as complex as the full history space $\mathcal{H}$. Rather, the theorem makes
explicit the structural content that any faithful completion must possess in order to support
counterfactual prediction under the intervention class $\mathcal{I}$.

\subsection{Memory and directedness as geometric consequences}

Within $\Sigma_2$, ``memory'' is not a hidden variable but a property of the fibered continuation
structure. When $\Pi$ is degenerate, a single reduced state $r$ corresponds to multiple histories
in $\mathcal{F}(r)$ reflecting different past interactions or correlations. If these fiber elements
support different admissible outcomes under the same intervention $I$, then reduced evolution is
necessarily history-dependent: the effective response depends on which fiber element is realized.
This dependence is what is operationally labeled ``memory'' in open systems.

``Directedness'' arises when forward interventions progressively constrain the feasible fiber.
Even if the underlying admissible-history structure on $\mathcal{H}$ is time-reversal symmetric,
the projection $\Pi$ can induce asymmetric continuation structure in $\mathcal{R}$: certain
intervention sequences reduce feasible subsets of $\mathcal{F}(r)$ in ways that cannot be undone
by any other interventions in $\mathcal{I}$. This geometric asymmetry appears as arrow-like
behavior at the reduced level without being postulated as a fundamental dynamical asymmetry.

Subsequent sections connect this structure to process-theoretic and causal-model formalisms and
discuss practical approximation schemes for working with $\Sigma_2$ objects.
\section{Relation to Process Tensors, Causal Models, Contextuality, and Measurement}
\label{sec:relations}

This work does not propose new dynamics or a replacement formalism for quantum theory. It instead
provides a structural diagnostic: it identifies when many-to-one operational reductions are
incapable—by construction—of supporting counterfactual reasoning under a specified intervention
class, and it characterizes the minimal representational structure required to restore such
reasoning.

\subsection{Process tensors and quantum combs}

Multi-time process tensors and quantum combs provide operational descriptions of non-Markovian
processes by encoding multi-time correlations into algebraic objects. In the present classification,
these are refined $\Sigma_1$ descriptions: they extend beyond instantaneous states but remain
\emph{algebraically closed}, in the sense that prediction proceeds by contraction or composition
rules internal to a fixed tensorial object.

The structural result of Section~3 applies equally to such descriptions. If the operational cut
identifying a process object is degenerate with respect to correlations that can later be probed
by interventions, then static algebraic closure does not guarantee faithfulness. In such cases,
distinct admissible histories may correspond to the same process tensor while supporting different
counterfactual continuations.

\paragraph{Schematic failure mode.}
Consider two admissible histories $h_1,h_2\in\mathcal{H}$ that induce the same static process tensor
$\mathcal{T}$ over a fixed set of times, so that all observable statistics for insertions at those
times coincide. Suppose, however, that $h_1$ and $h_2$ differ in correlations with degrees of freedom
discarded by the reduction defining $\mathcal{T}$. If the intervention class $\mathcal{I}$ includes
a later operation (e.g., free evolution or a control pulse) whose effect depends on those discarded
correlations, then the same $\mathcal{T}$ supports distinct admissible future outcomes depending on
which history was realized.

In this situation, the process tensor correctly predicts outcomes for the interventions it was
constructed to represent, but it cannot support counterfactual prediction for the enlarged
intervention class. The failure is not dynamical but representational: the tensor is a many-to-one
image of $\mathcal{H}$ that discards continuation-relevant distinctions.

From the $\Sigma_2$ perspective, the appropriate completion is not a single enlarged tensor, but a
family of algebraic descriptions indexed by fiber elements or continuation-equivalence classes in
$\mathcal{F}(r)$. Equivalently, one may view this as a sheaf-like structure in which $\Sigma_1$
descriptions are valid only on patches of the fiber determined by the intervention context.

This perspective is compatible with existing results emphasizing that extension theorems for
stochastic or quantum processes can fail in the presence of active interventions. Here, that
failure is diagnosed as an instance of an unfaithful operational cut rather than as a limitation
of a particular tensor formalism.

\subsection{Quantum causal models and interventions}

Quantum causal models and process-matrix frameworks focus on admissibility constraints at the level
of operations and correlations, often seeking conditions under which observed statistics admit a
causal explanation. The present framework addresses a prior question: whether the reduced
descriptions assigned to nodes support intervention semantics at all.

If nodal reductions are $\Sigma_1$ and unfaithful with respect to the intervention class relevant to
a causal analysis, then no causal structure built atop those reduced objects can consistently
support counterfactual reasoning without incorporating fiber-level admissibility structure. In
this sense, $\Sigma_2$ supplies a diagnostic precondition for causal modeling in reduced operational
spaces, rather than a competing causal formalism.

\subsection{Contextuality as unfaithful reduction}

Preparation contextuality provides a canonical instance of an unfaithful cut: distinct preparation
procedures yield operationally equivalent reduced descriptions (the same $\rho$) yet differ in
their statistics under subsequent measurement interventions. In the present language, $\Pi$ is the
map from preparation histories to reduced states, and unfaithfulness is precisely the failure of
continuation invariance over the fiber $\Pi^{-1}(\rho)$.

From the $\Sigma_2$ perspective, the appropriate descriptive object for preparations is not $\rho$
alone, but the lifted structure $\Sigma_2(\rho)$, which retains the intervention-indexed
continuation geometry required to support counterfactual prediction. This reframes contextuality
as a geometric feature of the reduction rather than as a purely interpretive puzzle about what the
quantum state ``really is.''

\subsection{Measurement and event structure}

Measurement problems can likewise be viewed through the lens of unfaithful cuts. Operational
reductions that identify macroscopically distinct branch histories under a single reduced
description destroy the continuation distinctions relevant to counterfactual questions of the
form: ``what would occur if we measured observable $M$ rather than $N$?''

In this perspective, the issue is not primarily to find a modified reduced dynamics, but to
recognize that an unfaithful reduction cannot support event-level counterfactuals without
additional structure encoding branch- or history-level admissibility. $\Sigma_2$ provides a
minimal language for representing such structure without committing to a particular interpretive
ontology or modification of quantum dynamics.
\section{Practical Computation and Approximation of \texorpdfstring{$\Sigma_2$}{Sigma-2} Objects}
\label{sec:sigma2-practical}

A natural objection is that $\Sigma_2$ objects appear computationally intractable: specifying
$\mathcal{F}(r)=\Pi^{-1}(r)$ and the relations $\{\mathcal{A}_I(r)\}$ exactly may be as hard as
working with $\mathcal{H}$ itself. The intent of $\Sigma_2$ is not to demand full reconstruction of
$\mathcal{H}$, but to identify the \emph{structurally necessary} information for counterfactual
prediction and to provide a target for approximation schemes.

\subsection{Finite parameterizations}

In many physically relevant cases, histories within $\mathcal{F}(r)$ differ only through a
finite set of parameters (e.g., a low-rank family of initial correlations, a small number of
bath modes, or a finite memory length). In such settings, one may represent $\mathcal{F}(r)$
implicitly as a parameterized family
\begin{equation}
\mathcal{F}(r) \approx \{ h(\theta) : \theta \in \Theta_r \},
\end{equation}
and approximate admissibility relations by pushforward of $\sim_I$ through this parameterization.

\subsection{Sample-based approximations}

When explicit parameterization is unavailable, $\Sigma_2$ admits sample-based approximations. One may
maintain a finite ensemble $\{h_j\}_{j=1}^N \subset \mathcal{F}(r)$ and approximate continuation
sets by evaluating which reduced outcomes are reachable under interventions for sampled elements.
This yields empirical approximations of the form
\begin{equation}
\widehat{\mathrm{Cont}}(r;\mathcal{I})
=
\{(I,\Pi(h_j')) : h_j \sim_I h_j' \text{ for some } j \}.
\end{equation}
Such approximations are directly relevant to robust control and verification problems, where the
objective is often to guarantee behavior over a set of possible hidden correlations rather than
to infer a unique hidden state.

\subsection{Coarse-graining by continuation equivalence}

Exact fibers are often too large, but $\Sigma_2$ suggests a principled coarse-graining: partition
$\mathcal{F}(r)$ into equivalence classes whose elements share identical continuation structure
under a restricted intervention class. Define an equivalence relation on $\mathcal{F}(r)$ by
\begin{equation}
h_1 \equiv_{\mathcal{I}} h_2
\iff
\forall I\in\mathcal{I}:\;
\{r':(h_1,r')\in\mathcal{A}_I(r)\}
=
\{r':(h_2,r')\in\mathcal{A}_I(r)\}.
\end{equation}
A $\Sigma_2$ description may then be represented by the quotient space $\mathcal{F}(r)/{\equiv_{\mathcal{I}}}$
together with the induced admissibility structure. This yields the \emph{coarsest representation
that preserves counterfactual prediction for the intervention class of interest}.

\subsection{Tractability and Structural Compression}
\label{sec:tractability-compression}

It is important to distinguish a $\Sigma_2$ description from a full
system--environment simulation. Although the fiber
$\mathcal{F}(r)=\Pi^{-1}(r)$ is defined as a set of histories, $\Sigma_2$
does not require enumerating or evolving those histories individually.
Rather, it tracks equivalence classes of histories that share identical
continuation structure with respect to the intervention class
$\mathcal{I}$.

In many physical regimes, the admissibility relations
$\{\mathcal{A}_I(r)\}_{I\in\mathcal{I}}$ collapse large families of
microscopically distinct histories into a low-dimensional
parameterization. In such cases, $\Sigma_2$ is strictly coarser than the
full history space $\mathcal{H}$ while remaining sufficient for
counterfactual prediction.

Examples of compressible fibers include:
\begin{itemize}
  \item finite-dimensional system--environment correlations, where the
  fiber is parameterized by a small set of correlation coefficients;
  \item bounded-memory reservoirs, where history dependence truncates
  after a finite time $\tau$;
  \item coarse-grained control-relevant observables, for which
  microscopically distinct states are indistinguishable to the available
  control bandwidth;
  \item experimentally indistinguishable preparation classes.
\end{itemize}

\subsection{The role of \texorpdfstring{$\Sigma_2$}{Sigma-2} in practice}

Even when full $\Sigma_2$ specification is impractical, the framework provides concrete benefits:
\begin{enumerate}
\item It supplies a diagnostic for when $\Sigma_1$ modeling is structurally incapable of supporting the
      counterfactual questions one intends to ask.
\item It defines the correct target object for approximation: not a single state or tensor, but a
      continuation-structured family indexed by hidden correlation classes.
\item It clarifies what must be preserved by reduced models and coarse-graining procedures to remain
      operationally faithful with respect to a chosen intervention class.
\end{enumerate}
\section{Applications and Connections}
\label{sec:applications}

This work does not propose new dynamics, a new interpretation, or a replacement formalism for
quantum theory. It provides a structural diagnostic: it identifies when an operational reduction
that discards causally relevant information cannot support counterfactual prediction under a
given intervention class, and it characterizes the minimal \emph{representational structure}
required to restore such predictive capacity. We now connect the $\Sigma_1/\Sigma_2$
classification to several concrete and active problem areas.

\subsection{Non-Markovian open quantum systems}
\label{sec:oqs}

Open quantum systems (OQS) theory typically seeks an effective evolution equation for a reduced
state $\rho_S(t)$, e.g.\ Lindblad or memory-kernel master equations, under assumptions such as
initial factorization or weak coupling. In our classification, these are $\Sigma_1$ descriptions:
prediction proceeds by applying an algebraically closed update rule on a reduced object.

The structural result of Section~3 identifies a standard mechanism by which $\Sigma_1$
descriptions can fail: whenever the reduction $\Pi$ identifies distinct correlated histories
as the same reduced state and the intervention class includes free evolution (or more generally,
a sufficiently rich set of admissible post-reduction operations), there need not exist a
single-valued rule on $\rho_S$ that supports counterfactual prediction across all admissible
histories.

This reframes persistent difficulties in strongly coupled or initially correlated regimes as a
\emph{structural} limitation of the reduction: the many-to-one nature of $\Pi$ induces a
degeneracy in continuation structure that cannot be repaired by improving the functional form
of a master equation on $\mathcal{R}$ alone. Classic discussions of how initial correlations
obstruct completely positive reduced dynamics can be viewed as early instances of this
unfaithfulness diagnosis \cite{Pechukas1994,Alicki1995}.

\subsection{Quantum control under memory and uncertain preparation}
\label{sec:control}

Quantum control practice often assumes that a reduced state (or an estimated effective process
model) suffices to predict responses to subsequent control pulses. In correlated or drifting
environments, however, identical $\Sigma_1$ reduced descriptions can respond differently to the
same control intervention. In our language, preparation ambiguity corresponds to fiber
degeneracy $\mathcal{F}(r)=\Pi^{-1}(r)$, and control success depends on which fiber element is
realized.

From the $\Sigma_2$ perspective, this suggests two complementary strategies:
\begin{enumerate}
  \item \textbf{Fiber localization:} augment experiments to infer which subset of $\mathcal{F}(r)$
        is realized (e.g.\ by probing correlation-sensitive observables).
  \item \textbf{Fiber-robust control:} design interventions $I$ such that the induced continuation
        set is effectively insensitive across the relevant equivalence classes of $\mathcal{F}(r)$,
        i.e.\ $\mathcal{A}_I(r)$ is approximately constant over the uncertainty region.
\end{enumerate}
This reframes the control design problem from \emph{reach a target state from $r$} to
\emph{achieve the desired outcome uniformly across a fiber or coarse-grained fiber class}.

\subsection{Preparation contextuality as an unfaithful cut}
\label{sec:contextuality}

Preparation contextuality formalizes the phenomenon that distinct preparation procedures can be
operationally equivalent at the level of a state description, yet yield different statistics
under subsequent measurement choices \cite{Spekkens2005}. Let $\Pi$ map preparation histories
to density matrices. Then preparation contextuality asserts the existence of $h_1,h_2$ with
$\Pi(h_1)=\Pi(h_2)=\rho$ but
$\mathrm{Cont}(\Pi(h_1);\mathcal{I})\neq\mathrm{Cont}(\Pi(h_2);\mathcal{I})$
for an intervention class $\mathcal{I}$ that includes measurement operations. This is precisely
the definition of an unfaithful cut.

From the $\Sigma_2$ perspective, the appropriate descriptive object for preparations is not
$\rho$ alone, but the lifted structure
$\Sigma_2(\rho)=(\rho,\mathcal{F}(\rho),\{\mathcal{A}_I(\rho)\}_{I\in\mathcal{I}})$,
which retains the intervention-indexed continuation geometry required for counterfactual
prediction. This does not re-derive contextuality theorems; it provides a unifying structural
interpretation in which contextuality appears as a special case of reduction-induced
unfaithfulness.

\subsection{Process tensors, quantum combs, and the ``static algebra'' limitation}
\label{sec:process-tensors}

Process tensors (and closely related quantum combs) provide an operational framework for
multi-time quantum processes by representing the dependence of outcomes on inserted operations
across multiple times \cite{Pollock2018PRA,Pollock2018PRL,Chiribella2009PRA,Chiribella2008PRL}.
These objects refine $\Sigma_1$ by encoding temporal correlations, while remaining algebraically
closed: prediction proceeds by contraction rules internal to a fixed tensorial object.

The present analysis emphasizes that static algebraic closure does not, by itself, guarantee
faithfulness. If the reduction identifying a ``process object'' is degenerate with respect to
causally relevant correlations not encoded in that object, then counterfactual interventions
can distinguish histories that the $\Sigma_1$ object treats as identical. In such cases, the
process object remains a many-to-one image of $\mathcal{H}$.

Related observations appear in work showing that extension theorems for stochastic or quantum
process descriptions can fail in the presence of active interventions \cite{Milz2020Quantum}.
The $\Sigma_2$ framework isolates the minimal geometric structure required to diagnose this
failure: rather than seeking a single ``complete'' tensor, one must track continuation
structure across fiber elements or coarse-grained fiber classes.

Reference~\cite{Milz2020Quantum} shows that quantum stochastic processes need not admit
Kolmogorov extensions when interventions are present: multi-time correlations may fail to be
consistently embeddable in a single static tensor structure. This can be understood as an
instance of unfaithful reduction: a process tensor constructed from limited observational data
may fail to support counterfactual predictions under interventions not included in its
construction. The $\Sigma_2$ framework diagnoses this failure as arising from fiber degeneracy in
the map from full histories to the process tensor and characterizes the minimal fiber-aware
augmentation required to restore predictive capacity.

\subsection{Quantum causal models and process-matrix frameworks}
\label{sec:causal}

Quantum causal models and the process-matrix framework study which correlations and operational
statistics admit causal explanation, including scenarios with indefinite causal order
\cite{Oreshkov2012,CostaShrapnel2016,WoodSpekkens2015}. These approaches typically operate at the
level of correlations among nodes and constraints on admissible process objects.

The present framework addresses a prior representational question: whether the reduced
descriptions assigned to nodes support intervention semantics at all. If nodal reductions are
unfaithful with respect to the intervention class relevant to a causal analysis, then no causal
structure built atop those reduced objects can yield valid counterfactual predictions without
incorporating fiber-level admissibility structure. In this sense, $\Sigma_2$ provides a
diagnostic precondition for causal modeling in reduced operational spaces.

\subsection{Measurement and branch structure (scope-limited)}
\label{sec:measurement}

The measurement problem is often framed as a tension between unitary evolution and definite
outcomes. We make a narrower claim: many measurement puzzles can be read as instances where a
reduction $\Pi$ identifies distinct histories (e.g.\ distinct macroscopic records) under a
single reduced description, thereby destroying the continuation structure required for
counterfactual prediction about measurement outcomes.

The $\Sigma_2$ framework does not select which history in a fiber is realized. It instead
clarifies what any successful account must retain: a branch- or fiber-indexed continuation
geometry sufficient to support counterfactual intervention questions of the form
``what would occur if we measured differently?'' This establishes a structural constraint on
reduction-based resolutions without committing to a particular ontology.

\subsection{Relation to classic no-go results}
\label{sec:nogos}

Classic results such as Kochen--Specker contextuality \cite{KochenSpecker1967} and Fine's
equivalence results for Bell-type models \cite{Fine1982} constrain classes of hidden-variable
assignments compatible with quantum statistics. The present work is orthogonal in emphasis:
rather than ruling out ontological models, it constrains the \emph{form of reduced descriptive
objects} under which counterfactual prediction is even well-defined. Where those no-go results
apply, they can often be reinterpreted as special instances of unfaithful cuts relative to
particular intervention classes.
\section{Practical Implementation and Approximation}
\label{sec:practical}

A natural objection is that $\Sigma_2$ objects appear to require enumerating large fibers
$F(r)=\Pi^{-1}(r)$ and intervention-indexed continuation structures. This section clarifies the
intended use: $\Sigma_2$ is not a promise of universal tractability; it is a specification of
\emph{structurally necessary information}, which enables principled approximation schemes.

\subsection{Finite parameterizations of fibers}
\label{sec:finite-param}

In many applications, the physically relevant variability within $F(r)$ is low-dimensional.
For example, when the dominant uncertainty is a small set of correlation parameters (e.g.\
effective bath-mode phases, coupling strengths, or a finite set of initial correlation
coefficients), one can represent a fiber by a parameter family
\[
  F(r) \approx \{ h(\theta) : \theta \in \Theta_r \subset \mathbb{R}^k \},
\]
and define continuation structure by propagating $\theta$-families under interventions.

\subsection{Sample-based approximation of continuation sets}
\label{sec:sampling}

When parametric structure is unknown or high-dimensional, one can approximate $\Sigma_2$ by
sampling histories consistent with a reduced description:
\[
  \widehat{F}(r) = \{ h^{(1)},\dots,h^{(N)} \} \subseteq \Pi^{-1}(r),
\]
then estimating continuation sets
\[
  \widehat{A}_I(r) = \{ \Pi(h') : \exists h\in \widehat{F}(r)\ \text{s.t.}\ h \sim_I h' \}.
\]
This yields empirical set-valued predictions: interventions map uncertainty in fibers to
uncertainty in outcomes.

\subsection{Coarse-graining by continuation equivalence}
\label{sec:coarse-grain}

A particularly useful reduction is to partition $F(r)$ into equivalence classes that share
continuation structure under a restricted intervention class $\mathcal{I}_0 \subseteq
\mathcal{I}$:
\[
  h \equiv_{\mathcal{I}_0} \tilde{h}
  \quad \Longleftrightarrow \quad
  \mathrm{Cont}(\Pi(h);\mathcal{I}_0)=\mathrm{Cont}(\Pi(\tilde{h});\mathcal{I}_0).
\]
The quotient $F(r)/\!\equiv_{\mathcal{I}_0}$ is a minimal coarse representation sufficient for
counterfactual prediction under $\mathcal{I}_0$.

\subsection{Detecting when \texorpdfstring{$\Sigma_1$}{Sigma-1} is adequate}
\label{sec:detect}

Although $\Sigma_1$ descriptions can fail generically - in the structural sense clarified above - under the conditions of Proposition~\ref{prop:unfaithful-from-correlations}, there are important regimes where it is adequate in
practice (e.g.\ engineered initial factorization, restricted intervention classes, weak coupling
and short times). Operationally, one can test for unfaithfulness by designing \emph{continuation
diagnostics}: prepare multiple realizations that agree in $r$ and apply a chosen intervention
family; statistically significant multi-modal outcome structure indicates fiber sensitivity.

\subsection{What \texorpdfstring{$\Sigma_2$}{Sigma-2} changes in practice}
\label{sec:practical-summary}

In applied settings, the principal role of $\Sigma_2$ is to:
\begin{itemize}
  \item identify the structural failure mode (degenerate reductions under correlation-sensitive
        interventions),
  \item provide a target object for approximation (fiber + continuation structure),
  \item support robust design goals (controls and inference procedures that remain valid over
        fibers or fiber classes).
\end{itemize}

\subsection{Experimental detection of unfaithful cuts}
\label{sec:experimental-unfaithful}

The unfaithful cut admits a direct experimental diagnostic. Consider a controlled quantum
system in which nominally identical reduced preparations $\rho_S$ may differ in correlations
with uncontrolled environmental degrees of freedom (e.g., due to preparation pathway
variations or ambient field fluctuations).

An experimental test proceeds as follows:
\begin{enumerate}
  \item Prepare an ensemble $\{h^{(1)}_k\}$ using preparation method $P_1$.
  \item Prepare an ensemble $\{h^{(2)}_k\}$ using preparation method $P_2$.
  \item Verify that both ensembles yield statistically indistinguishable reduced states
        $\rho_S$ under standard state tomography.
  \item Apply an identical control intervention $I \in \mathcal{I}$ to both ensembles.
  \item Measure the resulting outcome distributions.
\end{enumerate}

If the post-intervention outcome distributions differ beyond statistical uncertainty, this
demonstrates fiber degeneracy: the reduction $\Pi$ is unfaithful with respect to the intervention
$I$. Conversely, agreement of outcome statistics indicates that the experiment operates within
a $\Sigma_1$-adequate regime for the chosen intervention class.

This diagnostic distinguishes genuine second-order behavior from regimes in which preparation
variations are operationally irrelevant.
\section{Discussion and Outlook}
\label{sec:discussion}

We have shown that the common strategy of treating reduced algebraic objects (instantaneous
states, channels, or static multi-time tensors) as self-sufficient descriptions is structurally
limited. When the reduction map is degenerate in a way that discards causally relevant
correlations, and when the intervention class is rich enough to probe that dependence, the
resulting $\Sigma_1$ description is unfaithful: it cannot support counterfactual prediction
under any evolution rule defined solely on the reduced object. This limitation is
representational rather than dynamical and does not imply that $\Sigma_1$ descriptions are
universally invalid.

We then introduced $\Sigma_2$ geometry as an explicit representation of the minimal information
any faithful reduced description must encode: a base point $r$, its fiber
$\mathcal{F}(r)=\Pi^{-1}(r)$, and an intervention-indexed continuation structure
$\{\mathcal{A}_I(r)\}$, together with a set-valued update rule induced by interventions.
This construction restores counterfactual semantics by design and establishes minimality in a
\emph{representational} sense: any faithful completion must carry at least enough information to
reconstruct the $\Sigma_2$ continuation structure.

\subsection{What we are not claiming}

\begin{itemize}
  \item We do not claim that quantum theory is incorrect or incomplete as a predictive framework.
  \item We do not introduce new microscopic dynamics or modification rules.
  \item We do not propose hidden variables as an ontological commitment.
  \item We do not claim that $\Sigma_2$ descriptions are always computationally tractable in full generality.
\end{itemize}

\subsection{Open questions}

Several directions merit future work:
\begin{enumerate}
  \item \textbf{Categorical and sheaf-theoretic formulations.} The fiber/admissibility structure
        suggests natural connections to fibrations and sheaf-like semantics for local predictive
        patches.
  \item \textbf{Complexity of $\Sigma_2$ approximation.} Characterizing when low-dimensional or
        efficiently sampleable approximations exist.
  \item \textbf{Experimental diagnostics.} Designing continuation-based tests that isolate
        unfaithful cuts in controlled settings such as open quantum systems, quantum control,
        and contextuality experiments.
  \item \textbf{Intervention-class dependence.} Systematically classifying which intervention
        classes render a given reduction faithful or unfaithful.
\end{enumerate}

\subsection{A practical message}

When $\Sigma_1$ methods succeed, they do so in restricted regimes where discarded correlations
are dynamically irrelevant for the intervention class of interest. When they fail, the correct
response is not merely to search for a better reduced evolution equation, but to recognize that
the operational reduction has created a blind spot: counterfactual prediction requires tracking,
at minimum, fiber-level admissibility structure. $\Sigma_2$ provides a principled diagnostic
language for identifying this boundary and for guiding the construction of approximations that
respect it.
\appendix
\section{\texorpdfstring{Worked Example: $2$-dimensional system $\otimes$ $2$-dimensional environment}{Worked Example: 2D system x 2D environment}}
\label{app:worked-example}

This appendix gives a fully explicit instance of an unfaithful cut and shows how the associated
$\Sigma_2$ object is constructed.

\subsection{Setup}

Let $\mathcal{H}$ be the set of admissible histories of a joint system $SE$ with Hilbert space
$\mathcal{H}_S \otimes \mathcal{H}_E$, where $\dim \mathcal{H}_S=\dim \mathcal{H}_E=2$.
Let $\Pi$ be the reduction to the system density matrix at a chosen reduction event:
\[
  \Pi(h) = \rho_S(t_0) = \mathrm{Tr}_E\big[\rho_{SE}(t_0)\big].
\]
Let the intervention class $\mathcal{I}$ include free evolution for a fixed duration $\Delta t$,
implemented by a joint unitary $U_{SE}$:
\[
  I_{\Delta t}: \rho_{SE}(t_0) \mapsto \rho_{SE}(t_0+\Delta t) = U_{SE}\rho_{SE}(t_0)U_{SE}^\dagger,
\quad
  \rho_S(t_0+\Delta t)=\mathrm{Tr}_E[\rho_{SE}(t_0+\Delta t)].
\]
Define the intervention relation $h \sim_{I_{\Delta t}} h'$ to mean that $h'$ is a possible
post-intervention continuation of $h$ under $I_{\Delta t}$.

\subsection{Two distinct histories with identical reduced state}

Fix an orthonormal basis $\{|0\rangle_S,|1\rangle_S\}$ for $S$ and two distinct orthonormal
bases $\{|0\rangle_E,|1\rangle_E\}$ and $\{|0'\rangle_E,|1'\rangle_E\}$ for $E$.
Define two pure joint states at $t_0$:
\begin{align}
  |\psi_1\rangle_{SE} &= \sqrt{p}\,|0\rangle_S|0\rangle_E + \sqrt{1-p}\,|1\rangle_S|1\rangle_E,\\
  |\psi_2\rangle_{SE} &= \sqrt{p}\,|0\rangle_S|0'\rangle_E + \sqrt{1-p}\,|1\rangle_S|1'\rangle_E,
\end{align}
with $p\in(0,1)$ and with $\{|0\rangle_E,|1\rangle_E\}$ not equal to $\{|0'\rangle_E,|1'\rangle_E\}$.

Both reduce to the same system state:
\[
  \rho_S(t_0)
  = \mathrm{Tr}_E\big[|\psi_1\rangle\langle\psi_1|\big]
  = \mathrm{Tr}_E\big[|\psi_2\rangle\langle\psi_2|\big]
  = p\,|0\rangle\langle 0| + (1-p)\,|1\rangle\langle 1|.
\]
Thus, there exist distinct histories $h_1,h_2 \in \mathcal{H}$ with $\Pi(h_1)=\Pi(h_2)$.

\subsection{Explicit correlation-sensitive dynamics}

Choose an entangling unitary, e.g.\ controlled-NOT with $S$ as control and $E$ as target, in the
$\{|0\rangle_E,|1\rangle_E\}$ basis:
\[
  U_{SE} = |0\rangle\langle 0|_S \otimes \mathbb{I}_E + |1\rangle\langle 1|_S \otimes X_E,
\]
where $X_E|0\rangle_E=|1\rangle_E$ and $X_E|1\rangle_E=|0\rangle_E$.

Under this unitary, $|\psi_1\rangle_{SE}$ evolves to
\[
  |\psi_1'\rangle_{SE} = \sqrt{p}\,|0\rangle_S|0\rangle_E + \sqrt{1-p}\,|1\rangle_S X_E|1\rangle_E
  = \sqrt{p}\,|0\rangle_S|0\rangle_E + \sqrt{1-p}\,|1\rangle_S|0\rangle_E,
\]
which factorizes as $(\sqrt{p}\,|0\rangle_S+\sqrt{1-p}\,|1\rangle_S)\otimes|0\rangle_E$ and hence
yields a pure reduced state
\[
  \rho_S^{(1)}(t_0+\Delta t) = |\phi\rangle\langle\phi|,\quad |\phi\rangle=\sqrt{p}\,|0\rangle+\sqrt{1-p}\,|1\rangle.
\]

By contrast, if $|\psi_2\rangle_{SE}$ is defined using a different environment basis, then the
same $U_{SE}$ (defined relative to the unprimed basis) acts differently on the environment
components of $|\psi_2\rangle_{SE}$. For generic choices of $\{|0'\rangle,|1'\rangle\}$, the
resulting reduced state $\rho_S^{(2)}(t_0+\Delta t)$ differs from $|\phi\rangle\langle\phi|$.
Therefore,
\[
  \Pi(h_1)=\Pi(h_2)\quad\text{but}\quad
  \Pi(h_1') \neq \Pi(h_2')
\]
for $h_1 \sim_{I_{\Delta t}} h_1'$ and $h_2 \sim_{I_{\Delta t}} h_2'$.

This is a predictive failure: an agent who knows only $\rho_S(t_0)$ cannot predict the outcome
of applying the same free-evolution intervention $I_{\Delta t}$, because the result depends on
which fiber element (which correlation structure) is realized.

\subsection{Constructing the \texorpdfstring{$\Sigma_2$}{Sigma-2} object}

Let $r=\rho_S(t_0)$. The fiber is $F(r)=\Pi^{-1}(r)$, containing at least $h_1,h_2$.
For an intervention $I_{\Delta t}$, define the continuation set
\[
  A_{I_{\Delta t}}(r)
  = \Big\{ \Pi(h') \,\Big|\, \exists h\in \Pi^{-1}(r)\ \text{s.t.}\ h \sim_{I_{\Delta t}} h' \Big\}.
\]
In this example, $A_{I_{\Delta t}}(r)$ contains at least two distinct elements
$\rho_S^{(1)}(t_0+\Delta t)$ and $\rho_S^{(2)}(t_0+\Delta t)$, reflecting the set-valued nature of
faithful prediction under the unfaithful cut.

The associated $\Sigma_2$ description is
\[
  \Sigma_2(r) = \big(r,\ F(r),\ \{A_I(r)\}_{I\in\mathcal{I}}\big),
\]
which restores counterfactual semantics: given $I\in\mathcal{I}$, the correct prediction is not
a single output but the admissible outcome set determined by the fiber and the intervention
relation.

% ---------- Bibliography ----------
\printbibliography

\end{document}